The entry in row $i$ and column $j$ of a tableau $T \in \Syt(\lambda)$ is denoted by $T_{i,j}$. In addition, we define $\row_i(T) := \{T_{i,j} \mid 1 \le j \le \lambda_i\}$ as the set of entries in the $i$-th row of $T$. For example, if we consider the SYT shown in Figure~\ref{fig:SYT example}, then $T_{3,2} = 8$ and $\row_3(T) = \{5,8\}$.

\begin{theorem}[Gessel~\cite{knuth_classes_schur_positive}]
\label{thm:schur functions young tableaux}
    For every $\lambda \vdash N$, the set $\Syt(\lambda)$ is Schur-positive with respect to $\Des$. Moreover, $\Q(\Syt(\lambda)) = s_\lambda$.
\end{theorem}

In 2015, Adin and Roichman proved the following criterion.
\begin{theorem}[{\cite[Prop.\ 9.1]{adin_roichman_fine_set}}]
\label{thm:criterion schur positive young}
    A set $\A$ is symmetric with respect to $D:\A\to 2^{[N-1]}$ if and only if
    \[
        \sum_{a \in \A} \boldsymbol{t}^{D(a)} = \sum_{\lambda \vdash N} c_\lambda \sum_{T \in \Syt(\lambda)} \boldsymbol{t}^{\Des(T)}
    \]
    for some values $c_\lambda$, where $\boldsymbol{t}^J := \prod_{j \in J} t_j$ for $J \subseteq [N-1]$. The coefficients $c_\lambda$ are the Schur-coefficients of $\A$. Moreover, $\A$ is Schur-positive if and only if $c_{\lambda} \in \NN$ for all $\lambda \vdash N$.
\end{theorem}

This criterion implies that proving the Schur-positivity of a set is achievable by establishing a statistic-preserving bijection between the set and SYTs of shapes corresponding to a specific multiset.

In this paper, we will also apply a recently formulated criterion for symmetry~\cite{my_pattern_theorem}.
\begin{definition}
\label{def:respect composition}
    Let $\A$ be a finite set with a statistic $D:\A \to 2^{[N-1]}$. The set of elements that \emph{respect} a given composition $\alpha \vDash N$, denoted $\A_D(\alpha)$, consists of the elements $a \in \A$ such that $D(a) \subseteq S_\alpha$, where $S_\alpha$ is the set corresponding to the composition $\alpha$. When $D$ is clear from the context, we may write $\A(\alpha)$ instead.
\end{definition}

\begin{lemma}
\label{lem:my symmetry criterion}
    A set $\A$ is symmetric if and only if $|\A(\alpha)| = |\A(\beta)|$ for all $\alpha \sim \beta \vDash N$.
\end{lemma}

Note that only sets of permutations are considered in~\cite{my_pattern_theorem}. However, Lemma~\ref{lem:my symmetry criterion} applies to other sets as well.



    

    

Another useful result about symmetric sets and symmetric functions is due to Bloom and Sagan:
\begin{lemma}[Bloom and Sagan {\cite[Lemma 2.2]{BloomSagan}}]
\label{lem:singleton not symmetric}
    
    For every set $S \subseteq [N-1]$, the function $F_{N,S}$ is symmetric if and only if $S=[N-1]$ or $S=\emptyset$.
\end{lemma}


\section{Proof of Theorem~\ref{thm:criterion schur positive two rows}}
\label{sec:criterion two rows}

Denote $\odds(k) := \{1,3,5,\dots, 2k-1\}$ for a nonnegative integer $k$. In particular, $\odds(0) = \emptyset$.

Recall Definition~\ref{def:sparse} and Definition~\ref{def:elements with given stat}. We divide the assertions of Theorem~\ref{thm:criterion schur positive two rows} into two propositions:
\begin{lemma}
\label{lem:if schur positive two rows then sparse and good sizes}
    Let $\A$ be a symmetric set with respect to a statistic $D:\A \to 2^{[N-1]}$, and denote $n = \lfloor\frac{N}{2} \rfloor$. Assume that the Schur expansion of $\A$ is $\Q_{D}(\A) = \sum_{k=0}^n c_k s_{N-k,k}$ for some $c_k \in \ZZ$. Then
    $D$ is a sparse function, and for every sparse set $J \subseteq [N-1]$, the cardinality of the set $\A(D \supseteq J)$ depends on the size of $J$ only.
\end{lemma}
\begin{lemma}
\label{lem:if sparse and good sizes then schur positive two rows}
    Let $\A$ be a finite set with a sparse statistic $D:\A \to 2^{[N-1]}$, and denote $n = \lfloor\frac{N}{2} \rfloor$. Assume that for every sparse set $J \subseteq [N-1]$, the cardinality of the set $\A(D \supseteq J)$ depends on the size of $J$ only.
    Then $\A$ is Schur-positive, and its Schur expansion is
    \[
        \Q_{D} (\A) = \sum_{k=0}^n |\A(D=\odds(k))|\; s_{N-k,k}.
    \]
\end{lemma}

We prove Lemma~\ref{lem:if schur positive two rows then sparse and good sizes} in Section~\ref{subsec:proof lemma if schur positive two rows then sparse and good sizes}, and then show two proofs for Lemma~\ref{lem:if sparse and good sizes then schur positive two rows}. The first proof (presented in Section~\ref{subsec:proof if sparse and good sizes then schur positive two rows by induction}) is inductive. The second proof (presented in Section~\ref{subsec:proof if sparse and good sizes then schur positive two rows by superstandard}) is more involved, and it demonstrates the power of column superstandard tableaux, introduced by Hamaker, Pawlowski and Sagan~\cite{S3_patterns_list}, in proving Schur-positivity properties. We believe this approach can be applied in other cases as well.

Before proving these lemmas, let us prove a simple folklore lemma that will be useful for both lemmas:
\begin{lemma}
\label{lem:symmetric complement is symmetric}
    Let $\A$ be a symmetric set with respect to a statistic $D:\A \to 2^{[N-1]}$. Then $\A$ is also symmetric with respect to the complementary statistic $\bar{D}:\A \to 2^{[N-1]}$ defined by $a \mapsto [N-1]\setminus D(a)$.
\end{lemma}
\begin{proof}
    The set $\A$ is symmetric with respect to $D$, so by Theorem~\ref{thm:criterion schur positive young} there exist values $c_{\lambda},\ \lambda \vdash N$ such that
    \[
        \sum_{a \in \A} \boldsymbol{t}^{D(a)} = \sum_{\lambda \vdash N} c_\lambda \sum_{T \in \Syt(\lambda)} \boldsymbol{t}^{\Des(T)},
    \]
    where $\boldsymbol{t}^J := \prod_{j \in J} t_j$ for $J \subseteq [N-1]$, and therefore
    \[
        \sum_{a \in \A} \boldsymbol{t}^{[N-1]\setminus D(a)} = \sum_{\lambda \vdash N} c_\lambda \sum_{T \in \Syt(\lambda)} \boldsymbol{t}^{[N-1]\setminus \Des(T)} = \sum_{\lambda \vdash N} c_\lambda \sum_{T \in \Syt(\lambda')} \boldsymbol{t}^{\Des(T)},
    \]
    where $\lambda'$ is the partition conjugate to $\lambda$.
\end{proof}

\subsection{Proof of Lemma~\ref{lem:if schur positive two rows then sparse and good sizes}}
\label{subsec:proof lemma if schur positive two rows then sparse and good sizes}
\begin{proof}[\unskip\nopunct]
    The set $\A$ is symmetric, and its Schur expansion is $\Q_{D}(\A) = \sum_{k=0}^n c_k s_{N-k,k}$. Therefore, by Theorem~\ref{thm:criterion schur positive young}, we have
    \begin{equation}\label{eqn:proof if schur positive two rows then sparse}
        \sum_{a \in \A} \boldsymbol{t}^{D(a)} = \sum_{k=0}^n c_k \sum_{T \in \Syt(N-k,k)} \boldsymbol{t}^{\Des(T)}.
    \end{equation}
    First of all, notice that the set $\Des(T)$ is sparse for all $T \in \Syt(N-k,k)$. Therefore, by Equation~\eqref{eqn:proof if schur positive two rows then sparse}, the set $D(a)$ is sparse for all $a \in \A$.
    
    It remains to show that for every sparse set $J \subseteq [N-1]$ of size $k$, we have $\left|\A(D \supseteq J)\right| = \left|\A(D \supseteq \odds(k))\right|$. Denote $\alpha = (2^k,1^{N-2k}) \vDash N$. Additionally, let $i_1 < \dots < i_{N-k-1} \in [N-1]\setminus J$ denote the elements not in $J$. Since $J$ is sparse, the differences $i_{t+1} - i_t$ for $1 \le t \le N-k-2$, as well as $i_1$ and $N - i_{N-k-1}$, are either 1 or 2. Those that are equal to 2 correspond to the elements of $J$. The composition associated to the set $[N-1]\setminus J$ is denoted by $\beta := (i_1, i_2 - i_1, i_3 -i_2, \dots, i_{N-k-1} - i_{N-k-2}, N-i_{N-k-1}) \vDash N$. By Lemma~\ref{lem:symmetric complement is symmetric}, $\A$ is symmetric with respect to the complementary statistic $\bar{D}$. Notably, the compositions $\alpha$ and $\beta$ are equivalent, as they both have $k$ occurrences of $2$ and $N-2k$ occurrences of $1$. By Lemma~\ref{lem:my symmetry criterion}, we obtain that $|\A_{\bar{D}}(\beta)| = |\A_{\bar{D}}(\alpha)|$. By Definition~\ref{def:respect composition}, we obtain that $\A_{\bar{D}}(\beta) = \A(D \supseteq J)$ and $\A_{\bar{D}}(\alpha) = \A(D \supseteq \odds(k))$. Consequently, $|\A(D \supseteq J)| = |\A(D \supseteq \odds(k))|$, as required.
\end{proof}

\subsection{First proof of Lemma~\ref{lem:if sparse and good sizes then schur positive two rows}}
\label{subsec:proof if sparse and good sizes then schur positive two rows by induction}
We start by introducing an observation, which will serve as a crucial step in our inductive argument.
\begin{lemma}
\label{lem:remove first descent from syt}
    Let $k_1 \ge k_2$ be two nonnegative integers, and let $J \subseteq [k_1 + k_2 - 1]$ be a set. Then
    \[
        \left|\Syt(k_1, k_2)(\Des \supseteq J)\right| = \left|\Syt\big(k_1+1, k_2+1\big)\big(\Des \supseteq (\{1\}\cup (J+2))\big)\right|,
    \]
    where $J+2 = \{j + 2 \mid j \in J\}$.
\end{lemma}
\begin{proof}
    The proof is bijective: We associate a given tableau $T \in \Syt(k_1, k_2)$ such that $\Des(T) \supseteq J$ with the tableau $T' \in \Syt(k_1+1, k_2+1)$ defined by $T'_{1,1} = 1,\; T'_{2,1} = 2$, and for all $i \in \{1,2\}$ and $j \in [k_i]$, $T'_{i,j+1} = T_{i,j}+2$ (recall that $T_{i,j}$ is the entry in row $i$ and column $j$ of $T$). Intuitively, we first increase every entry of $T$ by $2$. Then, we insert a new column with the letters $1,2$ on the left of $T$, shifting the other columns to the right. Clearly, this process is injective, and $T'$ satisfies $T' \in \Syt(k_1+1, k_2+1)$ and $\Des(T) \supseteq \{1\}\cup (J+2)$. Furthermore, any $T' \in \Syt(k_1+1, k_2+1)$ with $\Des(T') \supseteq \{1\}\cup (J+2)$ can be obtained through this process.
\end{proof}

Now we are ready to prove Lemma~\ref{lem:if sparse and good sizes then schur positive two rows}.
\begin{proof}[First proof of Lemma~\ref{lem:if sparse and good sizes then schur positive two rows}]
    Let $\A$ be a finite set with a sparse statistic $D:\A \to 2^{[N-1]}$, and assume that for every sparse set $J \subseteq [N-1]$, the cardinality of the set $\A(D \supseteq J)$ depends on the size of $J$ only.
    By Theorem~\ref{thm:criterion schur positive young}, it suffices to show that
    \begin{equation}\label{eqn:first proof no consecutive indices is schur positive - generating function}
        \sum_{a \in \A} \boldsymbol{t}^{D(a)} = \sum_{k=0}^n |\A(D=\odds(k))| \sum_{T \in \Syt(N-k,k)} \boldsymbol{t}^{\Des(T)},
    \end{equation}
    where $n = \lfloor\frac{N}{2}\rfloor$.
    
    We prove it by induction on $N \in \NN$:
    
    For $N \le 1$, the statement holds trivially.
    
    Next, we assume that the statement holds for every value smaller than $N$ and proceed to prove it for $N$. Our goal is to establish the equation:
    \begin{equation}\label{eqn:first proof no consecutive indices is schur positive - set equality}
        \left|\A(D=J)\right| = \sum_{k=0}^n c_k \left|\Syt(N-k,k)(\Des=J)\right|
    \end{equation}
    for every set $J \subseteq [N-1]$, where $c_k = \left|\A(D=\odds(k))\right|$.
    
    We begin by considering the case $J = \emptyset$. In this case, the only SYT $T$ with $\Des(T) = \emptyset$ is the unique tableau of one row. Consequently, Equation~\eqref{eqn:first proof no consecutive indices is schur positive - set equality} simplifies to $\left|\A(D=\emptyset)\right| = \left|\A(D=\odds(0))\right|$, which obviously holds.
    
    It now remains to prove Equation~\eqref{eqn:first proof no consecutive indices is schur positive - set equality} for all $J \ne \emptyset$. By  the inclusion--exclusion principle, it suffices to prove that the following holds for all $J \ne \emptyset$:
    \begin{equation}\label{eqn:first proof no consecutive indices is schur positive - set containment}
        \left|\A(D \supseteq J)\right| = \sum_{k=0}^n c_k \left|\Syt(N-k,k)(\Des \supseteq J)\right|.
    \end{equation}
    
    The function $D$ is assumed to be sparse, and the set $\Des(T)$ is sparse for all $T \in \Syt(N-k,k)$. Therefore, if $J$ is not sparse then Equation~\eqref{eqn:first proof no consecutive indices is schur positive - set containment} holds.

    Next, let $J \ne \emptyset$ be a sparse set, and denote $|J| = i$. The set $\odds(i)$ is sparse too, and has $i$ elements. Therefore, by the assumptions of the lemma, we obtain that $\left|\A(D \supseteq J)\right| = \left|\A(D \supseteq \odds(i))\right|$.
    In addition, by Theorem~\ref{thm:schur functions young tableaux}, the set $\Syt(N-k,k)$ is Schur-positive and has the Schur expansion $\Q(\Syt(N-k,k)) = s_{N-k,k}$. Therefore, we can apply Lemma~\ref{lem:if schur positive two rows then sparse and good sizes} and deduce that
    \[
        \left|\Syt(N-k,k)(\Des \supseteq J)\right| = \left|\Syt(N-k,k)(\Des \supseteq \odds(i))\right|.
    \]
    Therefore, Equation~\ref{eqn:first proof no consecutive indices is schur positive - set containment} reduces to
    \begin{equation}\label{eqn:first proof no consecutive indices is schur positive - odds set}
        \left|\A(D \supseteq \odds(i))\right| = \sum_{k=0}^n c_k \left|\Syt(N-k,k)(\Des \supseteq \odds(i))\right|,
    \end{equation}
    where $i > 0$. Thus, the proof will be completed when we establish Equation~\eqref{eqn:first proof no consecutive indices is schur positive - odds set}.
    
    Denote $\A_1 = \{a \in \A \mid 1\in D(a)\}$. Furthermore, define a new statistic $D_1:\A_1 \to 2^{[N-3]}$ by $D_1(a) = (D(a)\setminus\{1\})-2$. Notice that $\A_1$ satisfies the requirements of Lemma~\ref{lem:if sparse and good sizes then schur positive two rows} with respect to $D_1$, where $N$ is replaced by $N-2$. Therefore, we may assume, by the induction hypothesis, that Lemma~\ref{lem:if sparse and good sizes then schur positive two rows} holds for $\A_1$. Thus, we can apply Equation~\eqref{eqn:first proof no consecutive indices is schur positive - set containment} for $\A_1$, while replacing $c_k$ with $|\A_1(D_1 = \odds(k))| = |\A(D = \odds(k+1))| = c_{k+1}$, and obtain
    \begin{equation}\label{eqn:first proof no consecutive indices is schur positive - induction}
        \left|\A_1(D_1 \supseteq J)\right| = \sum_{k=0}^{n-1} c_{k+1} \left|\Syt(N-2-k,k)(\Des \supseteq J)\right|.
    \end{equation}
    
    If we substitute $J = \odds(i-1)$ into Equation~\ref{eqn:first proof no consecutive indices is schur positive - induction}, we obtain
    \[
        \left|\A_1(D_1 \supseteq \odds(i-1))\right| = 
         \sum_{k=0}^{n-1} c_{k+1} \left|\Syt(N-2-k,k)(\Des \supseteq \odds(i-1))\right|.
    \]

    Notice that $\left|\A_1(D_1 \supseteq \odds(i-1))\right| = \left|\A(D \supseteq \odds(i))\right|$. Moreover, by Lemma~\ref{lem:remove first descent from syt}, the number of tableaux $T$ of shape $(N-k-1,k-1)$ with $\Des(T) \supseteq \odds(i-1)$ is equal to the number of tableaux $T$ of shape $(N-k,k)$ with $\Des(T) \supseteq \odds(i)$. Therefore, we can reformulate the equation and obtain that
    \[
        \left|\A(D \supseteq \odds(i))\right| = 
         \sum_{k=1}^{n} c_k \left|\Syt(N-k,k)(\Des \supseteq \odds(i))\right|.
    \]
    The only difference between this equation and Equation~\eqref{eqn:first proof no consecutive indices is schur positive - odds set} is the omission of the summand corresponding to $k=0$. However, this summand evaluates to zero and does not affect the overall expression. Indeed, $\left|\Syt(N,0)(\Des \supseteq \odds(i))\right| = 0$, since the unique SYT of shape $(N)$ has no descents.
\end{proof}


\subsection{Second proof of Lemma~\ref{lem:if sparse and good sizes then schur positive two rows}}
\label{subsec:proof if sparse and good sizes then schur positive two rows by superstandard}
As a first step of this proof, we prove that a set satisfying the conditions of Lemma~\ref{lem:if sparse and good sizes then schur positive two rows} is symmetric.
\begin{lemma}
\label{lem:sufficent condition symmetry two rows}
    Let $\A$ be a finite set with a sparse statistic $D:\A \to 2^{[N-1]}$, and denote $n = \lfloor\frac{N}{2} \rfloor$. Assume that there are constants $b_k \in \NN$ for $0 \le k \le n$, such that for every sparse set $J \subseteq [N-1]$, $|\A(D \supseteq J)| = b_{|J|}$.
    Then $\A$ is symmetric with respect to $D$.
\end{lemma}
\begin{proof}
    Following Lemma~\ref{lem:symmetric complement is symmetric}, it suffices to prove that $\A$ is symmetric with respect to the complementary statistic $\bar{D}$. As we proceed to prove it by Lemma~\ref{lem:my symmetry criterion}, let us find $|\A_{\bar{D}}(\alpha)|$ for all $\alpha = (\alpha_1, \dots, \alpha_\ell) \vDash N$.
    
    If a composition $\alpha$ has $\max_i (\alpha_i) > 2$, 
    then there exists $j \in [N-2]$ such that $j,j+1 \notin S_\alpha$, where $S_\alpha$ is the set corresponding to $\alpha$. By Definition~\ref{def:respect composition}, for every element $a \in \A_{\bar{D}}(\alpha)$ we have $\bar{D}(a) \subseteq S_\alpha$, so $j, j + 1 \notin \bar{D}(a)$. Consequently, $j,j+1 \in D(a)$, and the set $D(a)$ is not sparse. However,
    $D$ is assumed to be a sparse function, so $\A_{\bar{D}}(\alpha) = \emptyset$.
    
    Now assume that $\max_i (\alpha_i) \le 2$, and denote by $k = |\{i \mid \alpha_i = 2\}|$ the number of occurrences of $2$ in $\alpha$. Thus, we have $\ell = N - k$.
    Consider the set $J = [N-1] \setminus S_\alpha$. Notably, the set $J$ is sparse and consists of $k$ elements.
    Let $a \in \A$ be an element. We have $a \in \A_{\bar{D}}(\alpha)$ if and only if $\bar{D}(a) \subseteq S_\alpha$, or equivalently, $D(a) \supseteq J$. Thus, $|\A_{\bar{D}}(\alpha)| = |\A(D \supseteq J)| = b_k$.
    
    To conclude, if $\alpha$ contains an element larger than $2$ then $\A_{\bar{D}}(\alpha) = \emptyset$. Otherwise, the size of $\A_{\bar{D}}(\alpha)$ depends only on the number of occurrences of $2$ in $\alpha$. Therefore, $|\A_{\bar{D}}(\alpha)| = |\A_{\bar{D}}(\beta)|$ for all $\alpha \sim \beta \vDash N$. By Lemma~\ref{lem:my symmetry criterion}, we obtain that $\A$ is symmetric.
\end{proof}

As the next step of the proof, we prove that $\A$ is Schur-positive and find its Schur coefficients. For this, we define a total order on partitions $\lambda \vdash N$:
\begin{definition}
\label{def:conjugate order}
    Let $\lambda, \mu \vdash N$ be partitions of $N$. Let $\lambda_i' = |\{j \mid \lambda_j \ge i\}|$ denote the length of the $i$-th column in the Young diagram of $\lambda$. We say that $\mu$ is \emph{larger} than $\lambda$ in the \emph{conjugate order}, and denote $\mu' > \lambda'$, if there exists $i$ such that $\mu_j' = \lambda_j'$ for all $j < i$ and $\mu_i' > \lambda_i'$.
\end{definition}

Following Hamaker, Pawlowski and Sagan \cite[Section 5]{S3_patterns_list}, we define the \emph{column superstandard} Young tableau of shape $\lambda \vdash N$, obtained by filling the columns of the Young diagram of shape $\lambda$ one by one. We denote it by $T_\lambda \in \Syt(\lambda)$.
Formally, $(T_\lambda)_{i,j} = \lambda_1' + \cdots + \lambda_{j-1}' + i$. For example, if $\lambda = (4,2,2,1)$, then $T_\lambda$ is the SYT in Figure~\ref{fig:SYT_lambda example}.

